% Co-governed and enforced under the Sovereign Integrity Protocol License (SIP License v1.1):
% https://github.com/tensorrent/ACS-Framework/blob/main/LICENSE
\documentclass[11pt,a4paper]{article}

\usepackage{amsmath,amssymb,amsthm}
\usepackage[margin=1.05in]{geometry}
\usepackage[protrusion=true,expansion=false]{microtype}
\usepackage{booktabs}
\usepackage{longtable}
\usepackage{enumitem}
\usepackage{xcolor}
\usepackage{hyperref}
\hypersetup{colorlinks=true,linkcolor=blue!55!black,urlcolor=blue!60!black,citecolor=blue!55!black}

\newcommand{\dI}{\Delta\mathcal{I}}
\newcommand{\TT}{\textsf{T1}}
\newcommand{\Tw}{\textsf{T2}}
\newcommand{\Th}{\textsf{T3}}
\newcommand{\Tf}{\textsf{T4}}
\newcommand{\paper}[1]{\textsf{#1}}

\theoremstyle{definition}
\newtheorem{principle}{Principle}
\theoremstyle{plain}
\newtheorem{claim}{Claim}

\title{\textbf{One Mechanism, Many Forms}\\[0.45em]
\large The Information-Asymmetry Identity as a Single Generative Principle,\\
and Exactly How Far It Has Been Earned}
\author{Bradley Wallace\\
\small \textit{Independent Researcher / TensorRent}\\
\small with Claude (Anthropic) as instrument and aggregator\\
\small \texttt{One Mechanism Many Forms} (synthesis)}
\date{June 2026}

\begin{document}
\maketitle

\begin{abstract}
\noindent This paper states, and bounds, a single claim: the form/function asymmetry that recurs
across the ACS corpus is \emph{one mechanism expressed through many forms}, not a family of
analogues. The claim is ontological, not epistemic --- the forms are faithful instances of one
mechanism running, not lossy views of a hidden object. What makes ``one mechanism'' more than a
slogan is an \emph{identity chain}: the information-asymmetry functional $\dI$ equals the
transfer-entropy asymmetry (by definition), whose second-order Taylor coefficient equals the Lie
bracket $[f,g]$ (the BCH--TE morphism, \paper{FF06a}, \Tw), and which is identified with the
renormalization-group $c$-function (the central unifying conjecture, \Tw/\Th). The first two links
are established; the third is the load-bearing open identification, and we say so plainly. We then
catalogue the forms in which the mechanism is \emph{forced} --- gravity\,$\to$\,colour, the Riemann
zeros, the RG flow, the recursive inversion, the shuffle-knife separation, the prime-gap operator,
and the relational reversible-flattening case --- each cited to its source paper with a tier label.
Finally we draw the radius: the identity holds where it is proven and is a conjecture where it is
gestured (biology, cognition, institutions, language, and the theological reading), with each new
form a \emph{falsification test} rather than an assumption. Forms that looked like the mechanism and
failed are kept first-class. We do not prove the mechanism universal; we prove it singular across the
treated domains by identity, and we state the one link whose proof would settle the rest. Tiers:
\TT{} measured, \Tw{} proven/forced, \Th{} conjecture, \Tf{} falsified.
\end{abstract}

\tableofcontents

\section{The thesis, stated once}\label{sec:thesis}

The corpus has, until now, reported the same shape in many places without committing to what the
sameness \emph{is}. There are two ways to read a recurring shape, and they are not interchangeable.

The \emph{epistemic} reading: there is one hidden object, and each discipline --- physics,
number theory, computation, theology --- sees its own refraction of it through its own instrument.
This is the blind-men-and-the-elephant frame; the plurality lives in the observers, and every view
is a distortion of something unreachable. It is the wrong frame for this corpus, and it quietly
undoes the corpus's central result: \paper{FF06a} does not \emph{view} gravity and infer colour;
it shows that $\mathfrak{su}(3)$ \emph{is} the unique closure attractor of the Palatini bracket. A
forced result is not a refraction.

The \emph{ontological} reading, which this paper adopts: there is one mechanism, and it is
genuinely \emph{expressed} through many forms. The plurality lives in the substrates, and each form
is a faithful instance of the same mechanism running --- as the wave equation is not three refracted
glimpses of a hidden master-wave but the same differential structure instantiated in sound, light,
and quantum amplitude. Each instance is the whole mechanism, not a fragment of it.

\begin{principle}[One mechanism, many forms]
The information-asymmetry mechanism $\dI$ --- equivalently, the form/function coupling --- is a
single generative principle. Where the identity chain of Section~\ref{sec:identity} is established,
its appearances across domains are \emph{the same object in different substrates}, by identity, not
by analogy. Where the chain is only gestured, ``same mechanism'' is a conjecture whose test is
whether the chain instantiates there.
\end{principle}

The rest of the paper does three things and refuses to do a fourth. It states the identity chain
and labels each link's status (Section~\ref{sec:identity}); it catalogues the forms in which the
mechanism is forced (Section~\ref{sec:forms}); it draws the radius and lists the forms that failed
the test (Sections~\ref{sec:radius}--\ref{sec:negatives}). It does \emph{not} claim the mechanism
is universal; that claim is not earned, and saying it would be is the failure mode this corpus was
built to avoid.

\section{What makes it one: the identity chain}\label{sec:identity}

``Many forms'' is cheap; anything resembles anything at low enough resolution. ``One mechanism''
costs an identity. The mechanism is one because of a chain of equalities, of which two links are
established and one is the central open identification.

\paragraph{Link 1 --- $\dI$ is the transfer-entropy asymmetry (definitional).}
For two mutually constraining fields $f,g$, the net information one carries about the other,
$\dI(f,g)$, is the transfer-entropy asymmetry. This is a definition, not a result; it fixes the
object.

\paragraph{Link 2 --- the bracket is the second-order coefficient of $\dI$ (\paper{FF06a}, \Tw).}
The BCH--TE morphism (\paper{FF06a}, Lemma 2.10) establishes that the Lie bracket $[f,g]$ is the
exact second-order Taylor coefficient of $\dI(f,g)$. This is the link that lets an
information-theoretic asymmetry generate algebra: applied to the vierbein--connection pair it
produces $\mathfrak{sl}(4)$, and the rest of \paper{FF06a} follows. The link is a proven lemma.

\paragraph{Link 3 --- $\dI$ is the renormalization-group $c$-function (the central conjecture,
\Tw/\Th).}
The unifying identification is that $\dI$, monotone under the form\,$\to$\,function flow and zero at
fixed points, \emph{is} the Zamolodchikov $c$-function (2D) / Komargodski--Schwimmer $a$-function
(4D): form $=$ UV couplings, function $=$ IR couplings, the $\beta$-function $=$ the
form\,$\to$\,function coupling, and the $c$-theorem $=$ the monotonicity of $\dI$. The $c$- and
$a$-theorems are established mathematics; what is \emph{not} yet a theorem is the identification
$\dI \equiv c$. We therefore label Link 3 \Tw/\Th: argued from the shared monotonicity and
fixed-point structure, not proven as an isomorphism.

\begin{claim}[The chain is the unification]
If Link 3 is upgraded from identification to theorem, then ``one mechanism, many forms'' is a
theorem on every domain to which the $c$/$a$-function applies, because the three links compose into
a single object: $\dI = $ transfer-entropy asymmetry $=$ bracket-generator $=$ $c$-function. Until
then, the unification is exactly as strong as Link 3, and no stronger. This is the falsifiable spine
of the whole synthesis: if $\dI$ is \emph{not} the $c$-function, the corpus's cross-domain sameness
collapses from identity back to analogy.
\end{claim}

This is the honest centre of the paper. The grandeur of ``one mechanism'' is borrowed entirely from
one equals-sign whose proof is open. Everything downstream is conditional on it, and we keep the
condition visible.

\section{The forms in which the mechanism is forced}\label{sec:forms}

What earns the singular reading, short of Link 3, is that in each treated domain the mechanism is not
imposed but \emph{forced} --- the same form/function coupling, compelled by the structure rather than
chosen. We catalogue the forms with the tier of the forcing.

\begin{longtable}{p{0.30\textwidth} p{0.62\textwidth}}
\toprule
\textbf{Form (paper)} & \textbf{How the one mechanism is forced there} \\
\midrule
\endhead
Gravity $\to$ colour (\paper{FF06a}) & The Palatini bracket forces $\mathfrak{su}(3)$ as the
\emph{unique} 8-dimensional closure attractor ($D<10^{-14}$ over 50k samples); form (vierbein) and
function (connection) couple via $[e,\omega]$; $\gamma=0.274$, $\lambda=2\sqrt3/27$, three
generations from Jacobi truncation. \TT/\Tw \\
The Riemann zeros (\paper{FF06b}, \paper{FF06b$'$}) & The spectrum splits into form (GUE, generic)
and function (prime-dual, $\zeta$-specific); the transport obstruction's character is the
response-driven readout, forced non-abelian by the simplicity of $\mathfrak{sl}(4)$; central number
forced to the $B{-}L$ charge $1/3$. \TT/\Tw \\
The RG flow (Link 3) & Form $=$ UV couplings, function $=$ IR couplings, $\beta$ $=$ the coupling,
$c$-theorem $=$ monotone $\dI$. The $c$/$a$-theorems are proven; the identification of $\dI$ with $c$
is the open link. \Tw/\Th \\
The recursive inversion (\paper{FF06c}) & The mechanism applied to its own output: a solution carries
the asymmetry of its origin on its boundary and \emph{becomes} the next constraint
($C_0\to S_0\to C_1$, non-zero holonomy), with $\dI$ inverting. This is the form/function asymmetry
iterated. \TT/\Tw \\
The shuffle-knife separation (\paper{FF06e}) & The marginal-matched surrogate is the operator that
\emph{isolates} the mechanism's function-content from its form-content: arithmetic
$\sim$$11{,}500\sigma$ (function) versus spacing/counting/shape $0.4$--$1.0\sigma$ (form). \TT \\
The prime-gap operator (\paper{N3}) & $T_m=T_0+P_m$ on $(\mathbb{Z}/m\mathbb{Z})^{*}$: the kernel
is exactly the source (form) sector, $\dim\ker P_m=\varphi(m)$; the perturbation carries the function
content. \TT \\
The relational case (\paper{FF06K}, \paper{FF06h}/\paper{i}) & Reversible flattening $=$ injective
homomorphism into a pre-specified structure: homomorphism carries density (form-coupling),
injectivity carries identity (function); the seam is the non-injective convolution. \TT/\Th \\
The classical-gravity boundary (this corpus, dead-end) & The field is form (the multipole shadow it
carries) and loses function (the radial profile): a non-injective convolution seam. A boundary
instance --- and the reason \paper{FF06a} routes colour through the algebra (forced), not the field
(lossy). \TT \\
\bottomrule
\end{longtable}

\noindent The catalogue is the evidence for the singular reading at its current strength: not ``these
rhyme'' but ``in each, the same form/function coupling is compelled, and in the algebraic and
spectral cases it is the \emph{same} $\dI$ by Links 1--2.'' The cross-domain identity is exactly as
wide as Link 3 reaches; the forcing within each domain is what is independently established.

\section{The radius, and the falsification surface}\label{sec:radius}

A claim of universality that cannot fail is worthless; a claim of singularity by identity can fail,
form by form, and that is its value.

\paragraph{Where it is earned.} In the domains of Section~\ref{sec:forms}, the mechanism is forced
(\TT/\Tw), and in the algebraic/spectral cases it is the same $\dI$ by Links 1--2. There, ``one
mechanism'' is not an interpretation laid over the data; it is the data.

\paragraph{Where it is gestured.} \paper{FF06c} maps the inversion across physics, mathematics,
computation, institutions, cognition, and language; the present paper has \emph{not} instantiated the
identity chain in the last four. ``Same mechanism in cognition'' (or biology, or institutions, or the
theological reading) is therefore a conjecture, and a sharp one: it asserts that the identity chain
instantiates there. Its falsification surface is precise --- exhibit the domain, attempt the chain,
and either the form/function coupling is forced (the mechanism is there) or it merely resembles (a
false-fit wearing the mechanism's name).

\paragraph{The asymmetry that keeps it honest.} A mechanism universal \emph{by assumption} is
unfalsifiable and explains nothing; a mechanism singular \emph{by identity} is falsifiable in every
new substrate, because the identity chain either holds or does not. The strength of ``one mechanism,
many forms'' is exactly that it can be wrong about any given form --- and the corpus has been wrong
about several, on purpose (Section~\ref{sec:negatives}).

\section{Forms that failed the test (first-class)}\label{sec:negatives}

The radius is drawn as much by where the mechanism was expected and did not appear as by where it did.
Each entry is a form that resembled the mechanism and was rejected by test.

\begin{longtable}{p{0.40\textwidth} p{0.52\textwidth}}
\toprule
\textbf{Candidate form (\Tf)} & \textbf{Why it is not the mechanism} \\
\midrule
\endhead
$\varphi$ + inverse-silver ``structures noise'' & The combination is $25\times$ \emph{worse} than
$\varphi$ alone ($0.02132$ vs $0.00084$): it structures nothing. \\
Motif retrieval via gcd (\paper{FF06K}) & $55.15\%$ vs $83.13\%$ majority baseline --- below the
trivial predictor; the function-content is not gcd-shaped. \\
Difference-tone / Tartini off-diagonal (\paper{FF06b$'$}) & Consistent with an incommensurable null
($1.2\sigma$); the channel is empty by prime-log incommensurability (unique factorisation). \\
``Tension'' between GUE and arithmetic (this corpus) & False: the zeros carry both; they are not
mutually exclusive. \\
``Independence / cheap add-on'' of arithmetic to GUE & False: an additive build tops at peak
$\sim$$6$ while GUE-passing, versus $228$ for the zeros. \\
N3: uniform spectral contraction; character block-diagonalisation; universal renormalised law &
Three conjectured structures tested and rejected by the prime data. \\
$\tan\beta=1/2$; $\lambda$ as RG fixed point; WdW $=$ $\dI{=}0$ (quantum) (\paper{FF06a} corrections)
& Misidentifications corrected: $\tan\beta$ free; $\lambda$ algebraic (Koide); WdW demoted to
conjecture. \\
\bottomrule
\end{longtable}

\noindent These are not embarrassments to be hidden; they are the calibration of the claim. A
mechanism that fit all of them too would be a mechanism that fits anything, which is the same as
fitting nothing.

\section{Interpretive coda (bracketed)}\label{sec:coda}

The mechanism has reactional names. ``Nested self-referential emergence'' is the recursion of
Section~\ref{sec:forms} (\paper{FF06c}). ``Made in the image of the maker'' has a clean structural
translation --- \emph{the made carries the maker's generative form, which is why it can make in turn}
--- and that translation is the recursion again: the solution carries the asymmetry of its origin and
becomes the next constraint. ``Endless unfolding'' is the iteration, with the qualifier that the
record is unfolded-so-far with walls mapped (Sections~\ref{sec:radius}--\ref{sec:negatives}), not
proven endless. We record these as interpretive translations of the response-driven structure, not
as measured claims; they are how the one mechanism reads in older instruments, and the structural
version is the falsifiable one.

\section{Scope: what is and is not established}\label{sec:scope}

\paragraph{Established.} (i) The form/function mechanism is \emph{forced}, not imposed, in each
domain of Section~\ref{sec:forms} (\TT/\Tw). (ii) In the algebraic and spectral cases it is the same
$\dI$, by Links 1--2 of the identity chain (definition $+$ the BCH--TE lemma). (iii) The mechanism is
falsifiable form-by-form, and several candidate forms have been rejected (Section~\ref{sec:negatives}).

\paragraph{Not established.} (i) Link 3, the identification $\dI \equiv$ the RG $c$-function, is the
central open problem; the cross-domain \emph{identity} (as opposed to within-domain forcing) is
conditional on it. (ii) The mechanism is \emph{not} shown universal; cognition, biology, institutions,
language, and the theological reading remain conjecture, each a separate test. (iii) Nothing here
proves RH, constructs a Hilbert--P\'olya operator, or closes the $\sim$$1\%$ coupling residual; those
walls stand exactly where the source papers leave them.

\paragraph{The one sentence that survives compression.} \emph{It is one mechanism expressed through
many forms --- by identity where the chain is proven, by conjecture where it is gestured, and
falsifiable in every form either way.}

\vspace{1em}
\noindent\rule{\linewidth}{0.4pt}\\
\small This synthesis (\texttt{One Mechanism Many Forms}) sits over the result papers it cites:
\paper{FF06a} (Colour from Gravity), \paper{FF06b}/\paper{FF06b$'$} (the Riemann spectral side and the
transport obstruction), \paper{FF06c} (the Inversion Arc), \paper{FF06e} (the Shuffle Knife),
\paper{N3} (the prime-gap operator), \paper{LF01} (finite-window localisation), and \paper{FF06K} (the
reversible-flattening process record). It states what they share and, more carefully, exactly how far
the sharing has been earned. The unification is real to the radius of Link 1--2 and conjectural beyond
it; the honest measure of the program is that the boundary is drawn, not blurred.

\end{document}
